\section{主流 Actor-Critic 族与诊断}

\subsection{算法族与设计差异}

\begin{table}[H]
\centering
\begin{tabular}{p{0.23\textwidth}p{0.21\textwidth}p{0.23\textwidth}p{0.2\textwidth}}
\toprule
算法 & Critic 类型 & 策略更新 & 适合场景 \\
\midrule
PPO & V 的 soft target & clipped surrogate & 模型同步训练、CPU 集群 \\
A2C/A3C & 单步 V + entropy & 过采样与不同 worker & 齐心并行、多实例 \\
SAC & 双 Q + entropy max & off-policy replay buffer & 连续动作、需高样本效率 \\
DDPG/TD3 & Deterministic policy & deterministic actor gradient & 高维连续控制 \\
\bottomrule
\end{tabular}
\caption{不同 Actor-Critic 变体的核心差异}
\end{table}

\begin{knowledgebox}{从 Lecture slide 中抽象出来的设计维度}
讲者强调的差异主要来自三个维度：更新频率（on-policy vs off-policy），目标函数约束（clipping / KL / trust region），以及是否带 entropy 或自动调节 exploration。
\end{knowledgebox}

\begin{figure}[H]
\centering
\includegraphics[width=0.85\textwidth]{slide-12.jpg}
\caption{Lecture slides 中关于 PPO 与 SAC 的对比图，强调了 on-policy 的稳定性与 off-policy 的样本效率互补。}
\end{figure}

\subsection{训练监控与实验指标}

\begin{itemize}
    \item \textbf{KL divergence}：监控更新是否突破 trust region；
    \item \textbf{Advantage 分布}：过宽代表 high variance，过窄代表 bias 大；
    \item \textbf{Critic loss 曲线}：暴涨意味着 learning rate 太大或目标不稳定；
    \item \textbf{Entropy}：下降太快说明 exploration 落入 determinism。
\end{itemize}

\begin{warningbox}{指标之间的 tradeoff}
降低 critic loss 并不总是好事——只要 advantage 估计偏差大，Actor 就会收到错误信号。稳定训练需要同时看 loss、KL 和 policy entropy 的趋势。
\end{warningbox}

\subsection{本章小结}

主流 Actor-Critic 族在 trust region、entropy 和 off-policy buffer 上各有侧重；实战中需要观测多条曲线并配合 early stopping 或 checkpoint rollback。

